\documentclass{article}
\usepackage[margin=2cm]{geometry}
\usepackage{amsmath}
\usepackage{graphicx}
\usepackage{booktabs}
\usepackage[utf8]{inputenc}
\usepackage{ragged2e}
\setlength{\emergencystretch}{3em}
\usepackage{paracol}
\begin{document}
\sloppy

\begin{paracol}{2}
\RaggedRight

Proposition of Appendix shows that a sufficient condition
for the identifiability in the case of Gaussian and Boltzmann
linear policies is that the second moment matrix of the feature vector E\textasciitilde{}dz* [(s)(s)] is non-singular along with
the fact that the policy e plays each action with positive
probability for the Boltzmann policy.


Concentration ResultWe are now ready to present a concentration result, of independent interest, for the parameters
and the negative log-likelihood that represents the central
tool of our analysis (details and derivation in Appendix).
Under Assumption and Assumption, let D
\{si,ai)\}=1 be a dataset of n > O independent sam-
0argminocel(0) and 0* = argminecel(0) . If
the empirical FIM:


\begin{equation}
\begin{array} { r l } { \hat { \mathcal { F } } ( \theta ) = \frac { 1 } { n } \sum _ { i = 1 } ^ { n } \mathbb { E } _ { a \sim \pi _ { \theta } ( \cdot | s ) } \left [ \overline { \mathbf { t } } ( s, a, \theta ) \overline { \mathbf { t } } ( s, a, \theta ) ^ { T } \right ] \quad ( 1 ) } & { ( 1 ) } \\ & { \qquad }\end{array}
\end{equation}


has a positive minimum eigenvalue Amin > 0 for all 0 e O,
then, for any  e [0, 1], with probability at least 1  :


\begin{equation}
\begin{array} { r } { \cdot } \\ { \left \| \widehat { \boldsymbol { \theta } } - { \boldsymbol { \theta } } ^ { * } \right \| _ { 2 } \leqslant \frac { \sigma } { \hat { \lambda } _ { \operatorname* { m i n } } } \sqrt { \frac { 2 d } { n } \log \frac { 2 d } { \delta } }. } \end{array}
\end{equation}


Furthermore, with probability at least 1  , individually:


\begin{equation}
\begin{array} { r l } & { \quad \cdots \cdot _ { F } - \xi ( \theta ^ { * } ) \leqslant \frac { d ^ { 2 } \sigma ^ { 4 } } { \hat { \lambda } _ { \operatorname* { m i n } } ^ { 2 } n } \log \frac { 2 d } { \delta } } \\ & { \quad \widehat { \ell } ( \theta ^ { * * } ) - \widehat { \ell } ( \widehat { \theta } ) \leqslant \frac { d ^ { 2 } \sigma ^ { 4 } } { \hat { \lambda } _ { \operatorname* { m i n } } ^ { 2 } n } \log \frac { 2 d } { \delta }. } \\ & { \quad \cdot \quad \cdot \quad \cdot \quad}\end{array}
\end{equation}


The theorem shows that the L2norm of the difference between the maximum likelihood parameter 0 and the true parameter 0* concentrates with rate O(n-1/2) while the likelihood  and its expectation  concentrate with faster rate
O(n-1). Note that the result assumes that the empirical FIM
F(0) has a strictly positive eigenvalue Amin > 0. This condition can be enforced as long as the true Fisher matrix F(@)
has a positive minimum eigenvalue Amin, i.e. under identifiability assumption (Lemma) and given a sufficiently large
number of samples. Proposition of Appendix provides the
minimum number of samples such that with probability at
least 1   it holds that Amin > 0.


Identification Rule Analysis The goal of the analysis of
the identification rule is to find the critical value c(1) so that
the following probabilistic requirement is enforced.


Let 8 e [0, 1]. An identification rule producing I iscorrect if: Pr (I = I*)  8.
d*E[|\{i  I* : i E Ie\}|] the expected fraction of parameters that the agent does not control
selected by the identification rule and with  = d*E|\{i E
I* : i  Ic\}|] the expected fraction of parameters that the
agent does control not selected by the identification rule. We

\switchcolumn
\RaggedRight

now provide a result that bounds  and  and employs them
to derive correctness.


Let I be the set of parameter indexes selected by the
Identification Rule obtained using n > 0 i.i.d. samples
collected with g*, with 0* e O. Then, under Assumption and Assumption, let 0+ = arg minoce,l(0) for all
2
e0)- 10*)  c(1) it holds that:


\begin{equation}
\begin{array} { r l } & { \mathrm { e v } _ { i } ~ j \quad \mathrm { ~ : ~ " ~ } j > \mathrm { ~ : ~ c o s ~ m o n s ~ m a s } } \\ & { \alpha \leq 2 d \exp \left \{ - \frac { c ( 1 ) \lambda _ { \operatorname* { m i n } } ^ { 2 } n } { 1 6 d ^ { 2 } \sigma ^ { 4 } } \right \} } \\ & { \beta \leqslant \frac { 2 d - 1 } { d ^ { * } } \sum _ { i \in I ^ { * } } \exp \left \{ - \frac { ( l ( \theta _ { i } ^ { * } ) - l}}}\end{array}\right.
\end{equation}


Furthermore, the Identification Rule is (d  d* ) + d* )
correct.
Since  and  are functions of c(1), we could, in principle, employ Theorem to enforce a value , as in Definition, and derive c(1). However, Theorem is not very attractive in practice as it holds under an assumption regarding the
minimum eigenvalue of the FIM and the corresponding estimate, i.e. Amin 2w2, that cannot be verified in practice
22
since Amin is unknown. Similarly, the constants d*, 1(0*)
and 1(0*) are typically unknown. We will provide in Section
a heuristic for setting c(1).


\subsection*{Policy Space Identification in a Configurable
Environment}


The identification rules presented so far are unable to distinguish between a parameter set to zero because the agent
cannot control it, or because zero is its optimal value. To
overcome this issue, we employ the Conf-MDP properties
to select a configuration in which the parameters we want to
examine have an optimal value other than zero. Intuitively, if
we want to test whether the agent can control parameter 0,,
we should place the agent in an environment w, E  where
0, is maximally important for the optimal policy. This intuition is justified by Theorem, since to maximize the power
of the test (1 - ), all other things being equal, we should
maximize the log-likelihood gap 1(0*) - 1(0*), i.e. parameter 0, should be essential to justify the agent's behavior. Let
I e \{1, .., d\} be a set of parameter indexes we want to test,
our ideal goal is to find the environment w1 such that:


wE argmaxwen\{l0w-l0*w\}2
where 0*(w)EargmaxoeeJM(0) and 0(w) 
argmaxeee, JM(0) are the parameters of the optimal
policies in the environment M in IIo and Ho, respectively. Clearly, given the samples D collected with a sin-.
gle optimal policy *(wo) in a single environment Mwo,
solving problem (2) is hard as it requires performing an offdistribution optimization both on the space of policy parameters and configurations. For these reasons, we consider a
surrogate objective that assumes that the optimal parameter
in the new configuration can be reached by performing a sin-.
gle gradient step


Let I e \{1, .., d\} and I = \{1,..., d\}\I. For a vector v, we
denote with v| the vector obtained by setting to zero the

\end{paracol}

\end{document}
